\begin{frame}{Operaciones elementales por filas: intercambio}

La \alert{operación elemental tipo 1}
\[
\alert{F_i\leftrightarrow F_j}
\]
intercambia las filas \(i\) y \(j\):

\medskip

\[
\begin{pmatrix}
\vdots&&\vdots\\
a_{i1} & \cdots & a_{in}\\
\vdots &&\vdots\\
a_{j1} & \cdots & a_{jn}\\
\vdots &&\vdots
\end{pmatrix}
\quad
\xrightarrow{\;F_i\leftrightarrow F_j\;}
\quad
\begin{pmatrix}
\vdots &&\vdots\\
a_{j1} & \cdots & a_{jn}\\
\vdots &&\vdots\\
a_{i1} & \cdots & a_{in}\\
\vdots &&\vdots
\end{pmatrix}.
\]

\end{frame}


\begin{frame}{Operaciones elementales por filas: producto por un escalar}

La \alert{operación elemental tipo 2}
\[
\alert{\lambda F_i},
\qquad \lambda\neq0,
\]
multiplica la fila \(i\) por el escalar no nulo \(\lambda\in k\):

\medskip

\[
\begin{pmatrix}
\vdots &&\vdots\\
a_{i1} & \cdots & a_{in}\\
\vdots &&\vdots
\end{pmatrix}
\quad
\xrightarrow{\;\lambda F_i\;}
\quad
\begin{pmatrix}
\vdots &&\vdots\\
\lambda a_{i1} & \cdots & \lambda a_{in}\\
\vdots &&\vdots
\end{pmatrix}.
\]

\end{frame}

\begin{frame}{Operaciones elementales por filas: suma de un múltiplo}

La \alert{operación elemental tipo 3}
\[
\alert{F_i+\lambda F_j},
\qquad i\neq j,
\]
sustituye la fila \(i\) por \(F_i+\lambda F_j\), donde $\lambda\in k$ es un escalar:

\medskip

\[
\begin{pmatrix}
\vdots &&\vdots\\
a_{i1} & \cdots & a_{in}\\
\vdots &&\vdots\\
a_{j1} & \cdots & a_{jn}\\
\vdots &&\vdots
\end{pmatrix}
\quad
\xrightarrow{\;F_i+\lambda F_j\;}
\quad
\begin{pmatrix}
\vdots &&\vdots\\
a_{i1}+\lambda a_{j1}
&
\cdots
&
a_{in}+\lambda a_{jn}\\
\vdots &&\vdots\\
a_{j1} & \cdots & a_{jn}\\
\vdots &&\vdots
\end{pmatrix}.
\]

\end{frame}

\begin{frame}{Equivalencia por filas}

\begin{definition}
Sean \(A,B\in M_{m\times n}(k)\). Diremos que \(A\) y \(B\) son
\alert{equivalentes por filas}, y escribiremos
\[
A\sim B,
\]
si \(B\) se obtiene a partir de \(A\) mediante una sucesión finita de
operaciones elementales por filas.
\end{definition}

\end{frame}

\begin{frame}{Matrices escalonadas}

\begin{definition}
Sea \(A\in M_{m\times n}(k)\). El \alert{pivote} de una fila no nula
es su primera entrada no nula.

La matriz \(A\) se dice \alert{escalonada por filas} si:

\begin{enumerate}
    \item todas las filas nulas, si las hay, están situadas después
    de las filas no nulas;

    \item el pivote de cada fila no nula está situado estrictamente
    a la derecha del pivote de la fila anterior.
\end{enumerate}

\[
\begin{pmatrix}
0    &\alert{\ast} & \ast & \ast & \ast & \ast & \ast & \cdots & \ast \\
0    & 0    & 0    & \alert{\ast} & \ast & \ast & \ast & \cdots & \ast \\
0    & 0    & 0    & 0    & \alert{\ast} & \ast & \ast & \cdots & \ast \\
\vdots &&&&&&&&\vdots \\
0    & 0    & 0    & 0    & 0    & 0    & 0    & \cdots & 0
\end{pmatrix}.
\]

\end{definition}

\end{frame}

\begin{frame}{Una matriz que no está escalonada}

Consideremos la matriz
\[
A=
\begin{pmatrix}
\alert{2} & 1 & -1 & 3 & 0\\
0 & 0 & \alert{1} & 2 & 4\\
0 & \alert{3} & 0 & 1 & 2\\
0 & \alert{5} & 2 & 0 & 1
\end{pmatrix}.
\]

Los pivotes están situados, respectivamente, en las columnas
\[
1,\qquad 3,\qquad 2,\qquad 2.
\]

\medskip

La matriz \(A\) \alert{no está escalonada por filas} porque:

\begin{itemize}
    \item el pivote de la tercera fila está situado a la izquierda
    del pivote de la segunda fila;

    \item el pivote de la cuarta fila está situado debajo del pivote
    de la tercera fila.
\end{itemize}

Cualquiera de estas dos condiciones es suficiente para que \(A\) no esté escalonada.

\end{frame}

\begin{frame}[t,allowframebreaks]{Eliminación gaussiana}

\begin{theorem}
Toda matriz $A_{m\times n}$ es equivalente por filas a una matriz escalonada.
\end{theorem}

\begin{proof}
Si \(A=0\), ya es escalonada. Supongamos \(A\neq0\), y sea \(j\)
la primera columna de \(A\) que contiene alguna entrada no nula.

\[
\left(
\begin{array}{ccccccc}
0&\cdots&0&a_{1j}& * & \cdots & *\\
0&\cdots&0&a_{2j}& * & \cdots & *\\
\vdots&&\vdots&\vdots&\vdots&&\vdots\\
0&\cdots&0&a_{mj}& * & \cdots & *
\end{array}
\right).
\]

Mediante una operación elemental de tipo \(1\) (intercambio de filas), si es necesario,
podemos situar una entrada no nula \(a_{1j}\) en la primera fila.

\medskip

Mediante operaciones elementales de tipo \(3\)
anulamos todas las entradas situadas debajo de \(a_{1j}\).
\[
\left(
\begin{array}{ccccccc}
0&\cdots&0&a_{1j}& * & \cdots & *\\
0&\cdots&0&a_{2j}& * & \cdots & *\\
\vdots&&\vdots&\vdots&\vdots&&\vdots\\
0&\cdots&0&a_{mj}& * & \cdots & *
\end{array}
\right)\xrightarrow[i=2,\ldots,m]{F_i-\frac{a_{ij}}{a_{1j}}F_1}
\left(
\begin{array}{cccc|ccc}
0&\cdots&0&a_{1j}& * & \cdots & *\\
\hline
0&\cdots&0&0&\\
\vdots&&\vdots&\vdots&&B&\\
0&\cdots&0&0&
\end{array}
\right).
\]

Aplicamos el mismo procedimiento a la submatriz \(B\).

\medskip

Repitiendo sucesivamente:
\begin{enumerate}
    \item localizamos la primera columna no nula de la submatriz restante;
    \item mediante una operación de tipo \(1\), situamos una entrada
    no nula en su primera fila;
    \item mediante operaciones de tipo \(3\), anulamos todas las
    entradas situadas debajo de ella;
\end{enumerate}
y obtenemos, después de un número finito de pasos, una matriz escalonada.

\medskip

Cada paso consiste únicamente en operaciones elementales por filas;
por tanto, la matriz obtenida es equivalente por filas a \(A\).
\end{proof}

Este procedimiento se denomina \alert{método de eliminación gaussiana} o \alert{método de Gauss}.
\end{frame}

\begin{frame}[allowframebreaks]{Cálculo de una forma escalonada}

\[
\begin{pmatrix}
0&0&2&1&-1\\
2&4&2&0&8\\
0&0&0&2&-5\\
4&8&8&4&9
\end{pmatrix}
\xrightarrow{F_1\leftrightarrow F_2}
\begin{pmatrix}
2&4&2&0&8\\
0&0&2&1&-1\\
0&0&0&2&-5\\
4&8&8&4&9
\end{pmatrix}
\]

\[
\xrightarrow{F_4-2F_1}
\begin{pmatrix}
2&4&2&0&8\\
0&0&2&1&-1\\
0&0&0&2&-5\\
0&0&4&4&-7
\end{pmatrix}
\xrightarrow{F_4-2F_2}
\begin{pmatrix}
2&4&2&0&8\\
0&0&2&1&-1\\
0&0&0&2&-5\\
0&0&0&2&-5
\end{pmatrix}
\]

\[
\xrightarrow{F_4-F_3}
\begin{pmatrix}
2&4&2&0&8\\
0&0&2&1&-1\\
0&0&0&2&-5\\
0&0&0&0&0
\end{pmatrix}.
\]

La última matriz está escalonada por filas.

\end{frame}

\begin{frame}{Matrices escalonadas y reducidas}

\begin{definition}

Una matriz escalonada por filas se dice \alert{reducida} si, además:

\begin{enumerate}

    \item todos sus pivotes son iguales a \(1\);

    \item cada pivote es la única entrada no nula de su columna.
\end{enumerate}

\[
\begin{pmatrix}
0    & \alert{1}    & \ast & \alert{0}    & \alert{0}    & \ast & \ast & \cdots & \ast \\
0    & 0    & 0    & \alert{1}    & \alert{0}    & \ast & \ast & \cdots & \ast \\
0    & 0    & 0    & 0    & \alert{1}    & \ast & \ast & \cdots & \ast \\
\vdots &&&&&&&&\vdots \\
0    & 0    & 0    & 0    & 0    & 0    & 0    & \cdots & 0
\end{pmatrix}.
\]

\end{definition}

\end{frame}

\begin{frame}[allowframebreaks]{Forma escalonada reducida}

\begin{theorem}
Toda matriz $A$ es equivalente por filas a una matriz
escalonada y reducida.
\end{theorem}

\begin{proof}
Como toda matriz es equivalente
por filas a una escalonada, podemos suponer que $A$ ya lo es.
Multiplicando cada fila no nula de \(A\) por el inverso de su pivote,
mediante operaciones elementales de tipo \(2\), podemos suponer que
todos los pivotes son \(1\).

\[
\begin{pmatrix}
0    &\alert{a_{1j_1}} & \ast & \ast &\cdots & \ast & \ast &  \cdots \\
0    & 0    & 0    & \alert{a_{2j_2}} & \cdots &\ast & \ast & \cdots \\
\vdots &\vdots &\vdots &\vdots &&\vdots&\vdots \\
0    & 0    & 0    & 0 & \cdots & \alert{a_{ij_i}} & \ast & \cdots \\
\vdots &\vdots &\vdots &\vdots &&\vdots&\vdots
\end{pmatrix}
\xrightarrow{\frac{1}{a_{ij_i}}F_i}
\begin{pmatrix}
0    &\alert{1} & \ast & \ast & \cdots & \ast & \ast &  \cdots \\
0    & 0    & 0    & \alert{1} & \cdots & \ast & \ast & \cdots \\
\vdots &\vdots &\vdots &\vdots &&\vdots&\vdots\\
0    & 0    & 0    & 0 & \cdots   & \alert{1} & \ast & \cdots \\
\vdots &\vdots &\vdots &\vdots &&\vdots&\vdots
\end{pmatrix}.
\]

A continuación, empezando por el último pivote y procediendo hacia
arriba, mediante operaciones elementales de tipo \(3\) anulamos
todas las entradas situadas por encima de cada pivote.

\[
\begin{pmatrix}
0    &\alert{1} & \ast & \ast & \cdots & a_{1j_i} & \ast &  \cdots \\
0    & 0    & 0    & \alert{1} & \cdots & a_{2j_i} & \ast & \cdots \\
\vdots &\vdots &\vdots &\vdots &&\vdots&\vdots\\
0    & 0    & 0    & 0 & \cdots   & a_{kj_i} & \ast & \cdots \\
\vdots &\vdots &\vdots &\vdots &&\vdots&\vdots\\
0    & 0    & 0    & 0 & \cdots   & \alert{1} & \ast & \cdots \\
\vdots &\vdots &\vdots &\vdots &&\vdots&\vdots
\end{pmatrix}
\xrightarrow[k<i]{F_k-a_{kj_i}F_i}
\begin{pmatrix}
0    &\alert{1} & \ast & \ast & \cdots & \alert{0}& \ast &  \cdots \\
0    & 0    & 0    & \alert{1} & \cdots & \alert{0} & \ast & \cdots \\
\vdots &\vdots &\vdots &\vdots &&\vdots&\vdots\\
0    & 0    & 0    & 0 & \cdots   & \alert{0} & \ast & \cdots \\
\vdots &\vdots &\vdots &\vdots &&\vdots&\vdots\\
0    & 0    & 0    & 0 & \cdots   & \alert{1} & \ast & \cdots \\
\vdots &\vdots &\vdots &\vdots &&\vdots&\vdots
\end{pmatrix}.
\]

La matriz obtenida es escalonada y reducida, y equivalente por filas
a \(A\).
\end{proof}

\end{frame}

\begin{frame}[allowframebreaks]{De la forma escalonada a la reducida}

\[
\begin{pmatrix}
2&4&2&0&8\\
0&0&2&1&-1\\
0&0&0&2&-5\\
0&0&0&0&0
\end{pmatrix}
\xrightarrow{\substack{
\frac12F_1\\
\frac12F_2\\
\frac12F_3
}}
\begin{pmatrix}
1&2&1&0&4\\
0&0&1&\frac12&-\frac12\\
0&0&0&1&-\frac52\\
0&0&0&0&0
\end{pmatrix}
\]

\[
\xrightarrow{
F_2-\frac12F_3
}
\begin{pmatrix}
1&2&1&0&4\\
0&0&1&0&\frac34\\
0&0&0&1&-\frac52\\
0&0&0&0&0
\end{pmatrix}
\xrightarrow{
F_1-F_2
}
\begin{pmatrix}
1&2&0&0&\frac{13}{4}\\
0&0&1&0&\frac34\\
0&0&0&1&-\frac52\\
0&0&0&0&0
\end{pmatrix}.
\]

La última matriz está escalonada y reducida.

\end{frame}

\begin{frame}{Operaciones elementales inversas}

Toda operación elemental por filas puede \alert{deshacerse}
mediante otra operación elemental.

\medskip

\begin{enumerate}

\item Intercambiar dos filas:
\[
A
\xrightarrow{F_i\leftrightarrow F_j}
B
\xrightarrow{F_i\leftrightarrow F_j}
A.
\]

\item Multiplicar una fila por \(\lambda\neq0\):
\[
A
\xrightarrow{\lambda F_i}
B
\xrightarrow{\lambda^{-1}F_i}
A.
\]

\item Sumar a una fila un múltiplo de otra:
\[
A
\xrightarrow{F_i+\lambda F_j}
B
\xrightarrow{F_i-\lambda F_j}
A,
\qquad i\neq j.
\]

\end{enumerate}

\medskip

En cada caso, la segunda operación se denomina
\alert{operación elemental inversa} de la primera.

\end{frame}

\begin{frame}{Matrices elementales}

\begin{definition}
Una \alert{matriz elemental} es una matriz que se obtiene de la
matriz identidad \(I_m\) mediante una única operación elemental
por filas.
\end{definition}

\medskip

Hay, por tanto, tres tipos de matrices elementales:

\begin{enumerate}
    \item \(E_{ij}\): se obtiene intercambiando las filas \(i\) y \(j\)
    de \(I_m\).

    \item \(E_i(\lambda)\), \(\lambda\neq0\): se obtiene multiplicando
    la fila \(i\) de \(I_m\) por \(\lambda\).

    \item \(E_{ij}(\lambda)\), \(i\neq j\): se obtiene sumando a la
    fila \(i\) de \(I_m\) la fila \(j\) multiplicada por \(\lambda\).
\end{enumerate}

\end{frame}

\begin{frame}{Matrices elementales de tipo 1}

La matriz elemental correspondiente a
\(F_i\leftrightarrow F_j\) es

\[
E_{ij}=
\begin{pmatrix}
1      & 0      & \cdots & 0      & \cdots & 0      & \cdots & 0\\
0      & 1      & \ddots & \vdots &        & \vdots &        & \vdots\\
\vdots & \ddots & \ddots & 0      &        & 0      &        & \vdots\\
0      & \cdots & 0      & \alert{0} & \cdots & \alert{1} & \cdots & 0\\
\vdots &        &        & \vdots & \ddots & \vdots &        & \vdots\\
0      & \cdots & 0      & \alert{1} & \cdots & \alert{0} & \cdots & 0\\
\vdots &        &        & \vdots &        & \vdots & \ddots & 0\\
0      & \cdots & \cdots & 0      & \cdots & 0      & 0      & 1
\end{pmatrix}.
\]

\[
\hspace{34mm}
\underset{i}{\uparrow}
\hspace{19mm}
\underset{j}{\uparrow}
\]

Se obtiene de \(I_m\) intercambiando las filas \(i\) y \(j\).

\end{frame}

\begin{frame}{Matrices elementales de tipo 2}

La matriz elemental correspondiente a
\(\lambda F_i\), con \(\lambda\neq0\), es

\[
E_i(\lambda)=
\begin{pmatrix}
1      & 0      & \cdots & 0               & \cdots & 0\\
0      & 1      & \ddots & \vdots          &        & \vdots\\
\vdots & \ddots & \ddots & 0               &        & \vdots\\
0      & \cdots & 0      & \alert{\lambda} & \cdots & 0\\
\vdots &        &        & \vdots          & \ddots & 0\\
0      & \cdots & \cdots & 0               & 0      & 1
\end{pmatrix}.
\]

\[
\hspace{44mm}
\underset{i}{\uparrow}
\]

Se obtiene de \(I_m\) multiplicando la fila \(i\) por \(\lambda\).

\end{frame}

\begin{frame}{Matrices elementales de tipo 3}

La matriz elemental correspondiente a
\(F_i+\lambda F_j\), con \(i\neq j\), es

\[
E_{ij}(\lambda)=
\begin{pmatrix}
1      & 0      & \cdots & 0      & \cdots & 0               & \cdots & 0\\
0      & 1      & \ddots & \vdots &        & \vdots          &        & \vdots\\
\vdots & \ddots & \ddots & 0      &        & 0               &        & \vdots\\
0      & \cdots & 0      & 1      & \cdots & \alert{\lambda} & \cdots & 0\\
\vdots &        &        & \vdots & \ddots & \vdots          &        & \vdots\\
0      & \cdots & 0      & 0      & \cdots & 1               & \cdots & 0\\
\vdots &        &        & \vdots &        & \vdots          & \ddots & 0\\
0      & \cdots & \cdots & 0      & \cdots & 0               & 0      & 1
\end{pmatrix}.
\]

\[
\hspace{34mm}
\underset{i}{\uparrow}
\hspace{19mm}
\underset{j}{\uparrow}
\]

Se obtiene de \(I_m\) sumando a la fila \(i\) la fila \(j\)
multiplicada por \(\lambda\).

\end{frame}

\begin{frame}{Matrices elementales y operaciones por filas}

\begin{proposition}
Sea \(A\in M_{m\times n}(k)\). Entonces:

\begin{enumerate}

\item La operación de tipo \(1\) viene dada por
\[
A
\xrightarrow{F_i\leftrightarrow F_j}
E_{ij}A.
\]

\item La operación de tipo \(2\), con \(\lambda\neq0\), viene dada por
\[
A
\xrightarrow{\lambda F_i}
E_i(\lambda)A.
\]

\item La operación de tipo \(3\), con \(i\neq j\), viene dada por
\[
A
\xrightarrow{F_i+\lambda F_j}
E_{ij}(\lambda)A.
\]

\end{enumerate}
\end{proposition}


\end{frame}

\begin{frame}{Matrices invertibles}

\begin{definition}
Sea \(A\in M_{m\times n}(k)\). Diremos que \(A\) es
\alert{invertible} si existe una matriz \(B\in M_{n\times m}(k)\)
tal que
\[
AB=I_m
\qquad\text{y}\qquad
BA=I_n.
\]

En tal caso, \(B\) se denomina una \alert{inversa} de \(A\).
\end{definition}

\end{frame}

\begin{frame}{Unicidad de la inversa}

\begin{proposition}
Sea \(A\in M_{m\times n}(k)\) una matriz invertible. Entonces su
inversa es única.
\end{proposition}

\begin{proof}
Sean \(B,C\in M_{n\times m}(k)\) dos inversas de \(A\). Entonces
\[
B
=BI_m
=B(AC)
=(BA)C
=I_nC
=C.
\]
\end{proof}

Por tanto, podemos denotar la única \alert{inversa de \(A\)} por
\alert{\(A^{-1}.\)}

\end{frame}


\begin{frame}{Inversa de una matriz elemental}

\begin{proposition}
Toda matriz elemental es invertible. Más concretamente:

\begin{enumerate}
    \item \(E_{ij}^{-1}=E_{ij}\).

    \item
    \(E_i(\lambda)^{-1}=E_i(\lambda^{-1})\) para \(\lambda\neq0\).

    \item 
    \(E_{ij}(\lambda)^{-1}=E_{ij}(-\lambda)\) para \(i\neq j\).
\end{enumerate}
\end{proposition}

\begin{proof}
En cada caso, la matriz indicada como inversa corresponde a la
operación elemental inversa. Por tanto, al multiplicar ambas matrices
en cualquiera de los dos órdenes se obtiene \(I_m\).
\end{proof}

\end{frame}

\begin{frame}{Producto de matrices invertibles}

\begin{proposition}
Sean \(A\in M_{m\times n}(k)\) y \(B\in M_{n\times p}(k)\)
matrices invertibles. Entonces \(AB\) es invertible y
\[
(AB)^{-1}=B^{-1}A^{-1}.
\]
\end{proposition}
\begin{proof}
Por asociatividad del producto,
\[
(AB)(B^{-1}A^{-1})
=A(BB^{-1})A^{-1}
=AI_nA^{-1}
=AA^{-1}
=I_m.
\]
Por otra parte,
\[
(B^{-1}A^{-1})(AB)
=B^{-1}(A^{-1}A)B
=B^{-1}I_nB
=B^{-1}B
=I_p.
\]
Por tanto, \(AB\) es invertible y
\[
(AB)^{-1}=B^{-1}A^{-1}.\qedhere
\]
\end{proof}

\end{frame}

\begin{frame}{Forma escalonada y matrices elementales}

\begin{corollary}
Sea \(A\in M_{m\times n}(k)\) y sea \(E\) una forma escalonada por filas
de \(A\). Entonces existe una matriz invertible \(P\in M_m(k)\),
producto de matrices elementales, tal que
\[
E=PA.
\]
\end{corollary}

\begin{proof}
Como \(E\) se obtiene de \(A\) mediante una sucesión finita de
operaciones elementales por filas, existen matrices elementales
\[
E_1,\ldots,E_r\in M_m(k)
\]
tales que
\[
E=E_r\cdots E_1A.
\]

Tomando
\[
P=E_r\cdots E_1,
\]
obtenemos
\[
E=PA.
\]

Como cada \(E_i\) es invertible y el producto de matrices invertibles
es invertible, \(P\) es invertible.
\end{proof}

\end{frame}

\begin{frame}{Vectores}

\begin{definition}
Un \alert{vector} de \(k^n\) es una matriz columna de \(n\) entradas:
\[
\vec{v}=
\begin{pmatrix}
v_1\\
v_2\\
\vdots\\
v_n
\end{pmatrix},
\qquad v_1,\ldots,v_n\in k.
\]

Denotamos por
\[
k^n=M_{n\times 1}(k)
\]
el conjunto de todos los vectores de \(n\) entradas con coeficientes
en \(k\).
\end{definition}

\end{frame}

\begin{frame}{Sistemas de ecuaciones lineales}

\begin{definition}
Un \alert{sistema de \(m\) ecuaciones lineales con \(n\) incógnitas}
\(x_1,\ldots,x_n\) es una expresión de la forma
\[
\left\{
\begin{aligned}
a_{11}x_1+\cdots+a_{1n}x_n&=b_1,\\
a_{21}x_1+\cdots+a_{2n}x_n&=b_2,\\
&\ \vdots\\
a_{m1}x_1+\cdots+a_{mn}x_n&=b_m,
\end{aligned}
\right.
\]
donde \(a_{ij},b_i\in k\).

Una \alert{solución} del sistema es un vector
\[
x=
\begin{pmatrix}
x_1\\
\vdots\\
x_n
\end{pmatrix}\in k^n
\]
cuyas entradas satisfacen simultáneamente todas las ecuaciones.
\end{definition}

\end{frame}

\begin{frame}{Forma matricial de un sistema}

Consideremos el sistema
\[
\left\{
\begin{aligned}
a_{11}x_1+\cdots+a_{1n}x_n&=b_1,\\
&\ \vdots\\
a_{m1}x_1+\cdots+a_{mn}x_n&=b_m.
\end{aligned}
\right.
\]

Definimos
\[
A=
\begin{pmatrix}
a_{11}&\cdots&a_{1n}\\
\vdots&&\vdots\\
a_{m1}&\cdots&a_{mn}
\end{pmatrix},
\qquad
x=
\begin{pmatrix}
x_1\\
\vdots\\
x_n
\end{pmatrix},
\qquad
b=
\begin{pmatrix}
b_1\\
\vdots\\
b_m
\end{pmatrix}.
\]

Entonces el sistema puede escribirse en la forma
\[
\boxed{Ax=b}.
\]

En efecto, la \(i\)-ésima entrada de \(Ax\) es
\[
(Ax)_i=\sum_{j=1}^n a_{ij}x_j,
\]
que es precisamente el miembro izquierdo de la \(i\)-ésima ecuación.

\end{frame}

\begin{frame}{Matriz ampliada}

\begin{definition}
La \alert{matriz de coeficientes} del sistema
\[
Ax=b
\]
es la matriz \(A\).

Su \alert{matriz ampliada} es
\[
(A\mid b)=
\left(
\begin{array}{ccc|c}
a_{11}&\cdots&a_{1n}&b_1\\
\vdots&&\vdots&\vdots\\
a_{m1}&\cdots&a_{mn}&b_m
\end{array}
\right).
\]
\end{definition}

\medskip

Cada fila de \((A\mid b)\) representa una de las ecuaciones del sistema.

\end{frame}

\begin{frame}{Incógnitas pivote e incógnitas libres}

Supongamos que la matriz ampliada de un sistema está escalonada
y reducida.

\begin{definition}
Si un pivote está en la columna correspondiente a la incógnita
\(x_j\), diremos que \(x_j\) es una \alert{incógnita pivote}.

Las restantes incógnitas se denominan \alert{incógnitas libres}.
\end{definition}

\medskip

Por ejemplo, la matriz ampliada
\[
\left(
\begin{array}{ccccc|c}
1&2&0&-1&0&3\\
0&0&1& 4&0&2\\
0&0&0& 0&1&-1\\
0&0&0& 0&0&0
\end{array}
\right)
\]
tiene como incógnitas pivote
\[
x_1,\quad x_3,\quad x_5,
\]
y como incógnitas libres
\[
x_2,\quad x_4.
\]

\end{frame}

\begin{frame}{Soluciones de un sistema escalonado y reducido}

La matriz ampliada
\[
\left(
\begin{array}{ccccc|c}
1&2&0&-1&0&3\\
0&0&1& 4&0&2\\
0&0&0& 0&1&-1\\
0&0&0& 0&0&0
\end{array}
\right)
\]
representa el sistema
\[
\left\{
\begin{aligned}
x_1+2x_2-x_4&=3,\\
x_3+4x_4&=2,\\
x_5&=-1.
\end{aligned}
\right.
\]

Las incógnitas libres \(x_2,x_4\) pueden tomar valores arbitrarios.
Escribiendo
\[
x_2=\lambda,\qquad x_4=\mu,
\qquad \lambda,\mu\in k,
\]
las incógnitas pivote quedan determinadas:
\[
x_1=3-2\lambda+\mu,\qquad
x_3=2-4\mu,\qquad
x_5=-1.
\]

Por tanto, las soluciones son
\[
\boxed{
x=
\begin{pmatrix}
3-2\lambda+\mu\\
\lambda\\
2-4\mu\\
\mu\\
-1
\end{pmatrix},
\qquad \lambda,\mu\in k.
}
\]

\end{frame}